\documentclass[11pt,a4paper]{article}  % Déclare la classe du document.

\newcommand{\chaptermark}[2]{\og #2 \fg{} (#1)} % ligne uniquement pour éviter une erreur en cas d'article
\include{../1ere-preambule}

\usepackage{epic}   % extension et simplification de  l'environnement picture
\usepackage{graphicx}    % permet d'inclure des graphique externe dans un document
\usepackage{arydshln}   %permet d'avoir des lignes en pointillés

\newcommand{\lignes}[1]{
	\setlength{\unitlength}{1mm}
	\vskip 0,4cm\noindent
	\multido{\i=0+1}{#1}{%
		\noindent
		\begin{picture}(150,5)
			\linethickness{0,005mm}
			\put(-5,0){\dashline{1}(175,5)(5,5)}
		\end{picture}\vskip 0,001cm
	}
	\vskip 0cm
}

\rhead{préparation DS n$^{\circ}03$ - Automatismes}

\newcommand{\va}[1]{\lvert \ #1 \ \rvert}
\newcommand{\vaf}[1]{\left\lvert \ #1\  \right\rvert}

\begin{document}
	\graphicspath{{images/}}
	
%	\noindent Nom: ${\scriptstyle  .........................................................}$ \hskip 1.8cm Prénom: ${\scriptstyle  ...................................................}$ 
	
	\vspace{.5cm}
	
	\begin{center}
		
		
		\fbox{
			\begin{minipage}[]{15cm}
				\begin{center}
					\LARGE{\rule{0.5cm}{0pt}\rule[-0.25cm]{0pt}{0.9cm}
						Préparation DS n$^{\circ}03$ $-$ Automatismes\\
						Fraction, équation et inéquations 
					}
				\end{center}
			\end{minipage}
		}
	\end{center}

%\begin{center}\textbf{55min - Barème indicatif}\end{center}
%\begin{center}\textbf{Les élèves avec un tier-temps ne traitent pas les questions avec  le symbole \ding{46}}\end{center}
\begin{center}\textbf{Calculatrice non autorisée.\\Les résultats devront être justifiés (à l'aide de calcul ou de propriétés).}\end{center}



\exerciceDS{ $-$ équation du premier degré}

Résoudre dans $\mathbb{R}$ chacune des équations suivantes.
\begin{multicols}{3}
	\begin{enumMet}
		\item $12 x-3=4x+5$;
		\item $4 x-5=6 x+1$;
		
		\columnbreak
		\item $2x+8=5x-3$;
		\item $\dfrac{1}{4}x-3=\dfrac{1}{3}x+2$;
		
		\columnbreak
		\item $0,2 x-0,5=0,7x+0,3$;
		\item $\dfrac{1}{4} x-\dfrac{3}{5}=\dfrac{1}{3} x+\dfrac{3}{4}$;
	\end{enumMet}
\end{multicols}

\exerciceDS{ $-$ équation produit nul}

Résoudre dans $\mathbb{R}$ chacune des équations suivantes.
\begin{multicols}{3}
	\begin{enumMet}
		\item $4(x-6)(x-1)=0$;
		\item $5(2x-12)(5x-10)=0$;
		\item $-2(3 x+3)(-x+5)=0$.
	\end{enumMet}
\end{multicols}

\exerciceDS{ $-$ Résoudre une inéquation du premier degré}

Résoudre dans \R chacune des inéquations suivantes.
\begin{multicols}{3}
	\begin{enumMet}
		\item $-3 x +15 \geqslant 2x-5$;
		\item $-4 x+7 \leqslant 0$;
		\item $\dfrac{3}{4} x+\dfrac{1}{2} > \dfrac{1}{2}x+3$;
	\end{enumMet}
\end{multicols}

\exerciceDS{ $-$ équation de la forme $x^{2}=a$}

Résoudre dans \R les équations suivantes.
\begin{multicols}{4}
	\begin{enumMet}
		\item $x^{2}=4$;
		\item $x^{2}=5$;
		\item $2 x^{2}+3=0$;
		\item $x^{2}-7=0$.
	\end{enumMet}
\end{multicols}

\exerciceDS{ $-$ inéquation et tableau de signes}

\begin{enumMet}
	\item Résoudre dans $\mathbb{R}$ l'inéquation $4 x+1 \geqslant 0$.
	\item Résoudre dans $\mathbb{R}$ l'inéquation $-2 x+3 \geqslant 0$.
	\item A l'aide d'un tableau de signes Résoudre l'inéquation $(4 x+1)(-2 x+3)\geqslant0$.
\end{enumMet}


\end{document}

